\documentclass[a4paper,10pt]{article}
\usepackage{amsmath}
\usepackage{amssymb}
\usepackage[utf8]{inputenc}
\usepackage{enumitem} 
\usepackage{hyperref}
\usepackage{graphicx}

\title{M4C 2025 \--  Assignment 3} 
\author{Lan Stare}
\date{\today}

\begin{document}

\maketitle

% -------------------------------------------------------------
\subsection*{Problem 1}
 
\begin{enumerate}[label= (\alph*)]

 \item Euler's totient function $\phi(n)$ is the function that maps a natural number $n \in \mathbb{N}$ to its number of coprimes that are smaller than $n$. Because $n = p q$, $\phi(n)$ is the number of coprimes of $p q$ from 1 up to $p q$.
 \item Remark: I got a bit carried away with exercising my formal proving ability which means this question was answered more thoroughly than it was probably expected (which is becoming my trait at this point).
 
 Let's first prove the next statement: \textit{Intergers smaller than $n = p q$ are coprime to $p$ and $q$ iff they are coprime to $n$.}
 
 This can be explained by the definition of coprime numbers. Let's first prove the implication to the right. Let $m < n$ be a natural number that is coprime to $p$ and $q$. Two numbers $m$ and $n$ are coprime iff $\gcd(m, n) = 1$. Because $n = p q$, the common divisers of $m$ and $n$ can only be: $1, p, q, n$ by the definition of prime numbers. From this we can learn that $p | m \implies \gcd(m, n) \neq 1$ and symmetrically $q | m \implies \gcd(m, n) \neq 1$. Because $m$ is by assumption coprime to $p$ and $q$, it follows that $\gcd(m, n)$ can only be $1$ or $n$. But because $m < n$, it can only be 1.

 The implication to the left is easier. If we know that $\gcd(m, n) = 1$, we know that $p \not| m$ and $q \not| m$. Because $p$ and $q$ are prime, it naturally follows $m$ is coprime to both $p$ and $q$. $\square$

 From the statement above we can learn that the only two numbers that are not coprime to $n$ and are still smaller than $n$ are exactly $p$ and $q$ (so exactly 2 integers). If we don't limit ourselves to smaller coprime numbers, there are infinitely amount of coprimes to $n$ (which would be divisible by at least one of $p, q$ or $n$).
 
\item  Because $p$ and $q$ are primes, we know that all smaller numbers than them are their coprimes, so: $\varphi(p) = (p - 1)$ and $\varphi(q) = (q - 1)$. Because we proved that the coprimes to $n$ are exactly the primes to $p$ and the primes to $q$, we can derive the value of the totient function: \[\varphi(p q) = \varphi(p) \varphi(q) = (p - 1)(q - 1)\]

\end{enumerate}

% -------------------------------------------------------------

\subsection*{Problem 2}
 
\begin{enumerate}[label= (\alph*)]

 \item \ldots

\end{enumerate}

% ------------------------------------------------------------- 
\subsection*{Problem 3}
\begin{enumerate}[label= (\alph*)]

 \item \ldots

\end{enumerate}

% -------------------------------------------------------------
\subsection*{Problem 4}
 
\begin{enumerate}[label= (\alph*)]

 \item \ldots

\end{enumerate}

% -------------------------------------------------------------

\end{document}

